\begin{table}[H]
    \centering
    \scriptsize
    \caption{Perbandingan RMSE SARIMA dan STARIMA per Wilayah}
    \label{table:perbandingan-rmse-arima-starima}
    \renewcommand{\arraystretch}{1.25}

    \begin{tabular}{l l c}
        \toprule
        \textbf{Wilayah} & \textbf{Model} & \textbf{RMSE} \\
        \midrule

        \multirow{2}{*}{Barat}
            & SARIMA$(1,0,0) \times (0,1,0)_{12}$    & 5.082 \\
            & STARIMA$(0,0,0) \times (0,1,1_{0})_{12}$        & 2.962 \\
        \midrule

        \multirow{2}{*}{Selatan}
            & SARIMA$(1,0,0) \times (0,1,0)_{12}$    & 4.680 \\
            & STARIMA$(0,0,0) \times (0,1,1_{0})_{12}$   & 3.921 \\
        \midrule

        \multirow{2}{*}{Tengah}
            & SARIMA$(1,0,0) \times (0,1,0)_{12}$    & 4.434 \\
            & STARIMA$(0,0,0) \times (0,1,1_0)_{12}$        & 4.092 \\
        \midrule

        \multirow{2}{*}{Timur}
            & SARIMA$(1,0,0) \times (0,1,0)_{12}$    & 4.643 \\
            & STARIMA$(0,0,0) \times (0,1,1_0)_{12}$        & 1.331 \\
        \midrule

        \multirow{2}{*}{Utara}
            & SARIMA$(1,0,0) \times (0,1,0)_{12}$    & 5.016 \\
            & STARIMA$(0,0,0) \times (0,1,1_0)_{12}$        & 2.992 \\
        \bottomrule
    \end{tabular}
\end{table}
